\section{Testing}

Two entry points, both dependency-light and network-free:

\begin{verbatim}
python3 self_check.py
python3 tests/test_spectral_pipeline.py
\end{verbatim}

\code{self\_check.py} is a synthetic end-to-end smoke test of both engines
on the 358-mm reference configuration; it pins the scintillation-limited
photometric S/N ($\sim1.3\times10^3$) and the per-element spectroscopic S/N
(hundreds). \code{tests/test\_spectral\_pipeline.py} holds the unit-level
regressions: unit conversion of user curves, SVO photon/energy semantics,
hard coverage failures, FITS atmosphere layouts, and the v10 physics checks
--- Krisciunas--Schaefer scattering values, ecliptic symmetry and position
dependence of the sky, van Rhijn bounds, Pickering vs.\ $\sec z$,
parallactic-angle sign convention, Moffat wings and decentred coupling,
Filippenko dispersion scale, scintillation and quantization magnitudes,
SA100/SA200 dispersion, and extended-source area fractions. The file
inserts the repository root on \code{sys.path}, so it runs from any working
directory; it is also \code{pytest}-compatible (every test is a plain
\code{test\_*} function).

Rules: any bug fix gains a regression test that fails on the pre-fix code;
any new physical relation gains a test pinned to a published number, not to
the code's own output. The GUI is deliberately untested --- keep logic out
of it and in the engines, where it is testable headlessly.
